%\sukim{looks like it could be moved to the next section since it's a bit short as a section and still about our task.}

%\kenneth{Define the task and term here.}
A document $D$ contains $n$ figures, $F_1$ to $F_n$, where $F_i$ has a caption $C_i$ that was written by the document author.
In document $D$, $j$ sentences, $M_{i,1}$ to $M_{i,j}$, explicitly mention $F_i$ (\eg, ``As shown in $F_i$...'').
%The objective of this work is to automatically generate a high-quality caption $C_i'$ for figure $F_i$ based solely on all its mentions ($M_{i,1}$ to $M_{i,j}$.)
The objective of this work is to automatically generate a high-quality caption, $C_i'$, for figure $F_i$ using only its mentions ($M_{i,1}$ to $M_{i,j}$) and the surrounding text of the mentions in document $D$.
%We hypothesize that a high-quality caption for $F_i$, $C_i'$, can be automatically generated 
%\kenneth{Have a clear problem statement -- what do you expect or hope to find?}
%Automatic generation of high-quality captions for figures in a document $D$ with $n$ figures, $F_1$ to $F_n$, based solely on their mentions, $M_{i,1}$ to $M_{i,j}$, in the text of document $D$.
%The quality of the generated captions, $C_i'$, will be evaluated through two approaches: similarity to the original captions, $C_i$, and perception of target readers of document $D$.

%\paragraph{Terminology.}
%In the remaining of the paper, we used the following terminology:
%In the rest of the paper, 
%The following terms will be used to communicate based on the formulation above.
In the rest of the paper, we use these terms:







% Please add the following required packages to your document preamble:
% \usepackage{booktabs}
% \usepackage{multirow}
\begin{table*}[t]
    \centering \small
    \addtolength{\tabcolsep}{-1.1mm}
    %\vspace{-2mm}
    \begin{tabular}{@{}l@{\kern5pt}l@{\kern3pt}ccccccccccccccc@{}}
    \toprule
    \multicolumn{2}{@{}l}{\multirow{2}{*}{\textbf{Source}}} & \multicolumn{2}{c}{\textbf{Random}} & \multicolumn{2}{c}{\textbf{Mention}} & \multicolumn{2}{c}{\textbf{Paragraph}} & \textbf{OCR} & \multicolumn{2}{c}{\textbf{Window{[}0, 1{]}}} & \multicolumn{2}{c}{\textbf{Window{[}1, 0{]}}} & \multicolumn{2}{c}{\textbf{Window{[}1, 1{]}}} & \multicolumn{2}{c}{\textbf{Window{[}2, 2{]}}} \\ \cmidrule{3-17}
    \multicolumn{2}{c}{} & \textbf{S} & \textbf{P} & \textbf{-} & \textbf{+OCR} & \textbf{-} & \textbf{+OCR} & \textbf{-} & \textbf{-} & \textbf{+OCR} & \textbf{-} & \textbf{+OCR} & \textbf{-} & \textbf{+OCR} & \textbf{-} & \textbf{+OCR} \\ \midrule
    
    \multicolumn{2}{@{}l}{\textbf{Caption}} & 35.23 & 44.43 & 53.43 & 60.16 & 75.19 & \textbf{76.68} & 34.75 & 60.85 & 65.43 & 59.09 & 64.19 & 65.20 & 68.73 & 69.09 & 71.77 \\
    \multicolumn{2}{@{}l}{\textbf{Source}} & 32.52 & 19.52 & \textbf{39.51} & 18.78 & 12.53 & 9.39 & 20.79 & 30.49 & 17.19 & 32.40 & 17.33 & 25.10 & 15.55 & 19.84 & 13.45 \\ 
    \bottomrule
    \end{tabular}
    \addtolength{\tabcolsep}{+1.1mm}
    \caption{Macro coverage rates (percentage) between captions and relevant texts (S: Sentence and P: Paragraph). 
    Caption coverage gives the percentage of words in the caption that can be found in the source texts and vice versa (punctuation and stop words are excluded.)
    % S and P stand for Sentence, and Paragraph, respectively.
    The results show that 76.68\% of the words in captions could be found in Paragraph+OCR,
    % and 75.19\% could be found in Paragraph,
    motivating us to generate captions by text summarization.
    % (We calculated the percentage excluding punctuation and stop words.)
    }
    \label{tab:alignment-result}
    \vspace{-2mm}
\end{table*}

\begin{comment}

\begin{table}[t]
    \small
    \centering
    
    \caption{Macro coverage rates (percentage) between captions and models. Caption coverage gives the percentage of words in the caption that can be found in the source texts and vice versa. The results show that 76.68\% of the words in captions could be found in Paragraph+OCR.}
     \vspace{-2mm}    
    \label{tab:alignment-result}
    
    \begin{tabular}{lrr}
    \toprule 
    % \hline
    \multicolumn{1}{c}{\multirow{2}{*}{\textbf{Source Text}}} & \multicolumn{2}{c}{\textbf{Coverage}} \\ 
    \cmidrule{2-3}
     & \multicolumn{1}{c}{\textbf{Caption}} & \multicolumn{1}{c}{\textbf{Source Text}} \\ 
     \midrule
    %  \hline
    \textbf{Mention}                & 53.43  & \textbf{39.51}  \\
    \textbf{Mention+OCR}            & 60.16  & 18.78  \\
    \textbf{Paragraph}              & 75.19  & 12.53  \\
    \textbf{Paragraph+OCR}          & \textbf{76.68}  & 9.39   \\
    \textbf{OCR}                    & 34.75  & 20.79  \\ 
    % \hline
    \midrule
    \textbf{Random Sentence}         & 35.23  & 32.52  \\
    \textbf{Random Paragraph}       & 44.43  & 19.52  \\ 
    % \hline
    \midrule
    \textbf{Window [0, 1]}           & 60.85  & 30.49  \\
    \textbf{Window [0, 1]+OCR}       & 65.43  & 17.19  \\
    \textbf{Window [1, 0]}           & 59.09  & 32.40  \\
    \textbf{Window [1, 0]+OCR}       & 64.19  & 17.33  \\
    \textbf{Window [1, 1]}           & 65.20  & 25.10  \\
    \textbf{Window [1, 1]+OCR}       & 68.73  & 15.55  \\
    % \textbf{Window[0, 2]}             & 63.93  & 25.83  \\
    % \textbf{Window[0, 2]+OCR}         & 67.81  & 15.70  \\
    % \textbf{Window[2, 0]}             & 61.47  & 28.98  \\
    % \textbf{Window[2, 0]+OCR}         & 66.04  & 16.15  \\
    \textbf{Window [2, 2]}           & 69.09  & 19.84  \\
    \textbf{Window [2, 2]+OCR}       & 71.77  & 13.45  \\
    % \hline 
    \bottomrule
    \end{tabular}
    %\vspace{-1mm}
\end{table}

\end{comment}







\begin{comment}

In this paper, ``\textbf{Mention}'' refers to the sentence in an article that explicitly refers to the target figure, and 
``\textbf{Paragraph}'' refers to all the text under the same \texttt{<p>} tag (paragraph tag) as the Mention in the PDF parsing result.
As sentences near a Mention might be more relevant, we also extracted $n$ preceding sentences and $m$ following sentences to form \textbf{Windows~[n, m]}.
For example, ``Windows [1, 2]'' refers to a text snippet of five sentences, \ie, one preceding sentence, one Mention sentence, and two following sentences.

% \vspace{-1pc}

\end{comment}


\begin{itemize} [leftmargin=*,itemsep=0mm]

    \item 
    %A ``\textbf{Mention}'' refers to a sentence in a document that explicitly mentions the target figure, such as ``As shown in Figure 6 ...'' Since there might be multiple Mentions throughout the paper, without further specification, we referred to the first Mention.
    A \textbf{``Mention''} refers to a sentence in a document that explicitly mentions the target figure, \eg, ``As shown in Figure 6...''
    If there are multiple Mentions, the first Mention is referred to. %without further specification.




    
%\vspace{-.5pc}

    \item 
    
    %A ``\textbf{Paragraph}'' refers to all the text under the same \texttt{<p>} tag as the Mention in the PDF-parsing result.
    %A \textbf{``Paragraph''} refers to all the text under the same <p> tag as the Mention in the result of PDF parsing.
    A \textbf{``Paragraph''} refers to a section of text containing a Mention. In this work, the boundaries of a Paragraph are determined by the \texttt{<p>} tag produced by PDF parsing.









    
%\vspace{-1pc}

    \item 
    %As sentences near a Mention may also be relevant, we extracted $n$ preceding sentences and $m$ following sentences to form \textbf{Window[n, m]}. For example, ``Window[1, 2]'' refers to a text snippet of four sentences, \ie, one preceding sentence, one Mention sentence, and two following sentences.
    Sentences near a Mention may contain relevant information, so we extracted $n$ preceding sentences and $m$ following sentences to form the \textbf{``Window[n,~m]''} text snippet.
    For instance, ``Window[1,~2]'' refers to a snippet of four sentences, including one preceding sentence, the Mention sentence, and two following sentences. %\sukim{You define \textbf{``Window[n,~m]''} here, but it's not mentioned in the section 4 but appears in the experiment, which could make readers confusing. Since this is just another type of input, we could move it to the experiment or merge with \textbf{Paragraph} here.}

    %\cy{Should we also say something about OCR since it is used in the motivating analysis?}



    \item 
    An \textbf{``OCR''} refers to the textual information (\eg, legends, labels, etc.) extracted from the image, by optical character recognition (OCR) software.


\end{itemize}






